\subsection{李群的根基}

命题 23. 设 G 是有限维实或复李群，r 是 L(G) 的根基（第一章第 5 节定义 2），n 是 L(G) 的最大幂零理想（第一章第 4 节第 4 号）。设 R（分别地，N）是李代数为 r（分别地，n）的 G 的积分子群。则 R（分别地，N）是 G 的可解（分别地，幂零）李子群，并且在 G 的每个连续自同构下都不变。G 的每个连通可解（分别地，幂零）正规子群都包含于 R（分别地，N）中。

群 R 是可解的（第 6 节命题 19）。设 K = R。令 G' 为 G 的连通可解正规子群。则 \(\overline{G'}\) 是 G 的连通可解（第 1 号命题 1 的推论 2）正规李子群（第 8 节第 2 号定理 2）。因此 L \((\overline{G'})\) 是 L(G) 的可解理想，从而 L \((\overline{G'}) \subset r\)，且 \(\overline{G'} \subset R\)。特别地，\(\overline{R} \subset R\)，所以 R 是闭的，因而是 G 的李子群。设 K = C。令 H 为 G 的底层实李子群。若 \(r'\) 是 L(H) 的根基，则 \(i'r\) 是 L(H) 的可解理想，从而 \(r' = i'r'\)；于是 \(r \subset r' \subset r\)，由上面所述，R 在 H 中闭，因而在 G 中闭；所以 R 是 G 的李子群。G 的每个连通可解正规子群都是 H 的连通可解正规子群，因而包含于 R 中。因此我们已经证明，当 K = C 或 K = R 时，R 是 G 的最大连通可解正规子群；所以 R 在 G 的每个连续自同构下都不变。关于 N 的证明完全类似。

定义 1. 设 G 是有限维实或复李群。G 的根基是 G 的最大连通可解正规子群。

注。即使 G 是连通的，G 也可能有不包含于 G 的根基中的可解正规子群。

命题 24. 设 K = R 或 C。令 \(G_{1}\)、\(G_{2}\) 为两个有限维连通李群，\(R_{1}\) 和 \(R_{2}\) 为它们的根基，\(\phi\) 为从 \(G_{1}\) 到 \(G_{2}\) 的满射态射。则 \(\phi(R_{1}) = R_{2}\)。

由第 3 节第 8 号命题 28，\(\mathrm{L}(\phi)\) 是满射。因此 \(\mathrm{L}(\phi)(\mathrm{L}(\mathbb{R}_1)) = \mathrm{L}(\mathbb{R}_2)\)（第一章第 6 节命题 2 的推论 3）。令 \(i\) 为从 \(\mathbb{R}_1\) 到 \(\mathbb{G}_1\) 的典范嵌入。则 \(\phi \circ i\) 的像是 \(\mathbb{R}_2\)（第 6 节第 2 号命题 1 的推论 I）。

命题 25. 设 \(\mathbf{K} = \mathbf{R}\) 或 \(\mathbf{C}\)。令 \(\mathbf{G}_1, \mathbf{G}_2\) 为有限维连通李群，\(\mathbf{R}_1\) 和 \(\mathbf{R}_2\) 为它们的根基。\(\mathbf{G}_1 \times \mathbf{G}_2\) 的根基是 \(\mathbf{R}_1 \times \mathbf{R}_2\)。这由第一章第 5 节命题 4 推出。

